\begin{titlepage}
\begin{abstract}
The rapid integration of Large Language Models (LLMs) into command-and-control infrastructures necessitates rigorous, verifiable simulation environments to evaluate their strategic behaviors. Existing geopolitical platforms often rely on unstructured, text-based mechanics that restrict causal analysis and permit historical data leakage. This thesis introduces GeoMAS (Geopolitical Multi-Agent Simulation), a novel neuro-symbolic framework that establishes a \textit{Duality of Determinism} by grounding the stochastic, hierarchical reasoning of LLM-driven state cabinets within a mathematically rigid, procedurally generated environment. Through extensive experimental validation, we demonstrate that the system does not converge to chaotic warfare or a symmetric balance of power; instead, it self-organizes into a structurally bipolar configuration. Furthermore, by structurally decoupling private intent from public diplomatic declarations, the framework successfully quantifies strategic deception, revealing that global strategy dominates government type in predicting deceptive behavior and culminating in the emergent phenomenon of \textit{Moral Washing}. Crucially, when subjected to a global Pandemic shock, the system exhibits micro-economic resilience but macro-political fragility, resulting in an irreversible hegemonic rank reversal. Finally, through counterfactual injections and XAI envelope audits, we prove that macro-structural paradigm shifts can be causally traced back to precise cognitive friction and forced rationalization within a single agent's prompt. Ultimately, GeoMAS provides researchers and policymakers with a verifiable testbed to empirically measure, track, and causally audit the safety and reliability of autonomous AI agents before their deployment in high-stakes, adversarial environments.
\end{abstract}
\end{titlepage}